\documentclass[10pt,a4paper]{amsart}
\usepackage[margin=19mm]{geometry}
\usepackage{amsmath,amssymb,mathtools}
\usepackage[scr=rsfs,cal=euler]{mathalfa}
\usepackage{xfrac}
\usepackage[unicode,hidelinks,bookmarks=false]{hyperref}
\newcommand{\ie}{i.e.\ }
\newcommand{\cf}{cf.\ }
\newcommand{\resp}{respectively\ }
\newcommand{\Subsection}{\subsection}
\providecommand{\nobreakdash}{\nobreak\hbox{-}}
\setlength{\emergencystretch}{2em}
\multlinegap0pt
\usepackage{amssymb}
\usepackage[shortlabels]{enumitem}
\usepackage{tikz}
\newtheorem{lemma}{Lemma}
\title{Dirichlet Eigenvalue Approximation on Manifolds with Cylindrical Boundary}
\author{Anusha Bhattacharya}
\date{Source: arXiv:2603.12371v1 (CC BY 4.0)}
\begin{document}
\maketitle

\noindent\emph{Source and licence.} Dirichlet Eigenvalue Approximation on Manifolds with Cylindrical Boundary. Version arXiv:2603.12371v1 (CC BY 4.0), distributed by arXiv under the Creative Commons Attribution 4.0 International licence, http://creativecommons.org/licenses/by/4.0/, as recorded in the arXiv licence metadata for this exact version and shown on the abstract page https://arxiv.org/abs/2603.12371v1. Full licence notice: \texttt{licenses/arxiv-2603.12371-CC-BY-4.0.txt}.\par\smallskip

\noindent\textit{Editorial scope.} Counted unit: the article's lemma lem:discrete-green (source0.tex lines 247--255). The article's complete original proof of it is delivered with it (source0.tex lines 257--339, one \texttt{proof} environment of 83 lines). No mathematics is written, repaired or re-derived: the statement and the proof are the article's own lines, in the article's own notation, and no retained line is substituted or re-flowed.  No further source-local context is needed: every cross-reference of every retained line resolves inside the two delivered spans.  Nothing was omitted from any retained span, so the package carries no omission record.  No retained line contains a citation, so the wrapper carries no bibliography and no bibliography entry.  No retained line contains a figure environment, an image-inclusion command or drawing code, so no image is delivered and the package declares no asset dependency.  Editorial formatting only, in addition: an AMS \texttt{amsart} preamble that restates the article's own declarations and macros used by the retained lines, namely the environments \texttt{lemma} on separate counters, together with the standard packages those lines need.  The retained environments print in this document as Lemma 1 in order of appearance, whereas the article prints the same items with the numbering of its own full document.  The complete native source of the article as distributed in the arXiv e-print for this exact version is bundled unchanged as \texttt{sources/196272/source0.tex}.
\begin{lemma}
\label{lem:discrete-green}
For all $u,v \in L^2(X_t)$,
\[
\langle \delta u, \delta v \rangle_{L^2(E)}
=
-\langle \Delta_\Gamma u, v \rangle_{L^2(X_t)} .
\]
\end{lemma}

\begin{proof}

\begin{align*}
\langle \delta u, \delta u \rangle_{L^2(E)}
&=
\frac{1}{2}\sum_{x_i\sim x_j}
w_{ij}\bigl(u(x_j)-u(x_i)\bigr)^2 \\
&=
\frac{1}{2}\sum_{x_i\sim x_j}
w_{ij}\bigl(u(x_i)-u(x_j)\bigr)u(x_i)
-
\frac{1}{2}\sum_{x_i\sim x_j}
w_{ij}\bigl(u(x_i)-u(x_j)\bigr)u(x_j).
\end{align*}

Since the edges are directed, interchanging the indices $i$ and $j$ in the second sum and using the symmetry
$w_{ij}=w_{ji}$, we obtain
\[
\langle \delta u, \delta u \rangle_{L^2(E)}
=
\sum_{x_i}
\left(
\sum_{x_j\sim x_i}
w_{ij}\bigl(u(x_i)-u(x_j)\bigr)
\right)
u(x_i).
\]
Since $u\in L^2(X_t)$, $u(x_i)=0$, for all $x_i\notin X_t.$ Hence, 
                            \[
\langle \delta u, \delta u \rangle_{L^2(E)}
=
\sum_{x_i\in X_t}
\left(
\sum_{x_j\sim x_i}
w_{ij}\bigl(u(x_i)-u(x_j)\bigr)
\right)
u(x_i).
\]
Now for $x_i\in X_t$,
\[
\Delta_\Gamma u(x_i)
=
\frac{1}{\mu_i}
\sum_{x_j\sim x_i}
w_{ij}\bigl(u(x_j)-u(x_i)\bigr),
\]
which implies
\[
\sum_{x_j\sim x_i}
w_{ij}\bigl(u(x_i)-u(x_j)\bigr)
=
-\mu_i\,\Delta_\Gamma u(x_i).
\]

Substituting this identity yields
\[
\langle \delta u, \delta u \rangle_{L^2(E)}
=
-\sum_{x_i}
\mu_i\,\Delta_\Gamma u(x_i)\,u(x_i)
=
-\langle \Delta_\Gamma u, u \rangle_{L^2(X_t)}.
\]

Finally, both mappings
\[
(u,v)\mapsto \langle \delta u,\delta v\rangle_{L^2(E)},
\qquad
(u,v)\mapsto \langle \Delta_\Gamma u,v\rangle_{L^2(X_t)}
\]
define symmetric bilinear forms on $L^2(X_t)$.
Therefore, the general identity follows from the polarization formula,
\[
\langle \delta u,\delta v\rangle
=
\frac14\bigl(
\langle \delta(u+v),\delta(u+v)\rangle
-
\langle \delta(u-v),\delta(u-v)\rangle
\bigr).
\]

\end{proof}

\end{document}
